\documentclass[a4paper,12pt]{article}

\usepackage{graphicx}
\usepackage{svg}
\usepackage{geometry}
\usepackage{eso-pic}
\usepackage{tikz}
\usetikzlibrary{calc}

\geometry{margin=0cm}

%----------------------------------------------------
% Fondo SVG
%----------------------------------------------------
\AddToShipoutPictureBG*{%
    \AtPageCenter{%
        \makebox(0,0){%
            \includesvg[
                width=\paperwidth,
                height=\paperheight,
                keepaspectratio=false
            ]{dibujof}
        }%
    }
}

\pagestyle{empty}

\begin{document}

\null

%----------------------------------------------------
% Texto dentro del marco
%----------------------------------------------------
\begin{tikzpicture}[remember picture,overlay]

\node[
    anchor=north west,
    text width=9cm,
    align=justify,
    font=\footnotesize
]
at ($(current page.north west)+(5.8cm,-4.5cm)$)
{

\begin{center}
{\normalsize\texttt{About me}}
\end{center}

\texttt{
Soy Ingeniero en Física Industrial, apasionado de las ciencias exactas, los números y los análisis. En este blog comparto mis conocimientos, proyectos personales y reflexiones sobre matemáticas aplicadas, física, análisis cuantitativo y tecnología.
}

\vspace{0.25cm}

\begin{center}
\scalebox{0.26}{%
    \includesvg{dibujoe}
}
\end{center}

\vspace{0.25cm}

\texttt{
Mis habilidades incluyen análisis de datos, softwares estadísticos y lenguajes de programación numéricos. Me especializo en la aplicación de métodos cuantitativos para resolver problemas complejos en diferentes áreas, combinando la rigurosidad científica con herramientas computacionales avanzadas.
}

};

\end{tikzpicture}

\end{document}
